\subsection*{5.6.}
要对不带支撑的上同调得到与上述类似的结果, 似乎必须有奇点消解可用.
在维数至多为 $1$ 时, 推论~5.3 已解决此问题.

\subsubsection*{定理 5.6.}
设 $N$ 为整数, 并假定对 $\mathbb F_q$ 上维数至多为 $N$ 的有限型概形,
可以作 Hironaka 意义下的强奇点消解 (因而目前可以取 $N=2$~[1]).
设 $f:X\longrightarrow Y$ 是 $\mathbb F_q$ 上有限型概形的态射,
$\underline G$ 是 $X$ 上可构造的 $T$-整 $E$ 层.
若 $\dim X\leq N$, 则 $Y$ 上的 $E$ 层 $R^if_*\underline G$ 均为 $T$-整的.

\medskip
\noindent\emph{证明.}
对 $X$ 的维数 $n\leq N$ 归纳. 设 $j:U\hookrightarrow X$ 为 $X$ 的开子集,
光滑且仿射, 使 $\underline G|_U$ 为扭常值层, 且补集 $Z$ 的维数小于 $n$.
由奇点消解, 层 $R^ij_*(j^*\underline G)$ 可构造 (SGA~5, Expos\'e~I; [2]);
当 $i>0$ 时它们支撑于 $Z$, 当 $i=0$ 时 $j_*(j^*\underline G)$ 与 $\underline G$
仅在 $Z$ 上不同. 因此由归纳假设, 只需对层 $j^*\underline G$ 及态射 $j$ 和 $f\circ j$
证明定理. 所以还可设 $X$ 光滑仿射, $\underline G$ 扭常值,
即来自 $X$ 的基本群在有限维 $E$ 向量空间上的连续表示.

设 $R$ 是 $E$ 的整数环, $\lambda$ 是素元, $e$ 是绝对分歧指数,
并选取整数 $b>e/(\ell-1)$. 由于 $X$ 光滑从而正规, 且其基本群为概有限群,
存在扭常值 $R$ 层 $\underline G^\circ$, 使得
\[
\underline G\simeq\underline G^\circ\otimes_R E.
\]
设 $p:X'\longrightarrow X$ 为使局部常值层
$\underline G^\circ\otimes_R R/\lambda^b$ 平凡化的平展覆盖.
层 $\underline G$ 是 $p_*p^*\underline G$ 的直和项, 故只需对 $p^*\underline G$
及 $f\circ p$ 证明定理. 尚需处理满足下述条件的情形:
\[
(*)\quad
\begin{minipage}{0.82\linewidth}
$X$ 光滑仿射, 存在扭常值 $R$ 层 $\underline G^\circ$, 使
$\underline G\simeq\underline G^\circ\otimes_R E$, 且
$\underline G^\circ\otimes_R R/\lambda^b$ 为常值层.
\end{minipage}
\]

将 $f$ 分解为固有态射与开浸入的复合, $f=\bar f\circ j$:
\begin{center}
\begin{tikzcd}
X \arrow[r,hook,"j",swap] \arrow[d,"f"'] &
X' \arrow[dl,"\bar f"]\\
Y & {}
\end{tikzcd}
\qquad .
\end{center}
定理对 $\bar f$ 已证 (5.2), 故只需对 $j$ 证明.
对 $X'$ 作奇点消解后, 可设 $X'$ 光滑, 且 $X$ 是 $X'$ 中某正规交叉除子的补集,
该除子为光滑除子之并. 假设 $(*)$ 蕴涵 $\underline G$ 沿此除子温和分歧.
此时有:

\subsubsection*{引理 5.6.1.}
设 $X$ 是 $\mathbb F_q$ 上的光滑有限型概形, $D$ 是 $X$ 中的正规交叉除子,
为光滑除子 $D_i$ 之并, $j:U\hookrightarrow X$ 为 $D$ 在 $X$ 中的补集,
$\underline G$ 是 $U$ 上沿 $D$ 温和分歧的扭常值 $E$ 层.
若 $\underline G$ 是 $T$-整的, 则层 $R^ij_*\underline G$ 均为 $T$-整的.

我们承认下述结果 (参见 SGA~1, Expos\'e~XIII).

\subsubsection*{引理 5.6.2.}
\begin{enumerate}[label=(\roman*)]
\item 在引理~5.6.1 的假设下, 设 $D'$ 为 $D$ 的一个分支,
  $D''$ 为其余分支之并, 置 $U'=X-D''$ 及 $D^{\prime\circ}=D'\cap U'$:
  \begin{center}
  \begin{tikzcd}
  {} & D^{\prime\circ}
    \arrow[r,hook,"j'_D",swap] \arrow[d,"i^\circ"] &
    D' \arrow[d,"i"]\\
  U \arrow[r,hook,"j_1",swap] &
    U' \arrow[r,hook,"j'",swap] & X
  \end{tikzcd}
  \qquad .
  \end{center}
  则基变换映射
  \[
  i^*Rj'_*(Rj_{1*}\underline G)
  \longrightarrow
  Rj'_{D*}\,i^{\circ *}(Rj_{1*}\underline G)
  \]
  为同构, 且层 $i^{\circ *}Rj_{1*}\underline G$ 扭常值,
  并沿 $D'-D^{\prime\circ}$ 温和分歧.

\item 若 $C$ 是 $X$ 上的光滑曲线, 在 $D$ 的光滑点 $x$ 处与 $D$ 横截,
  \begin{center}
  \begin{tikzcd}
  C\cap U \arrow[r,"j_c",swap] \arrow[d,"k'"] &
    C \arrow[d,"k"]\\
  U \arrow[r,"j",swap] & X
  \end{tikzcd}
  \qquad ,
  \end{center}
  则基变换映射
  \[
  k^*Rj_*\underline G\longrightarrow Rj_{c*}\,k^{\prime *}\underline G
  \]
  在 $x$ 处为同构.
\end{enumerate}

现对维数归纳, 从引理~5.6.2 推出引理~5.6.1, 暂时承认曲线的引理~5.6.1.
引理~5.6.2(ii) 与此假设表明, 层 $R^ij_*\underline G$ 在 $D$ 的每个光滑点处均为 $T$-整的.
特别地, 采用引理~5.6.2(i) 的记号, 层 $R^ij_{1*}\underline G$ 均为 $T$-整的.
因此引理~5.6.2(i) 与用于 $i^{\circ *}Rj_{1*}\underline G$ 及 $j'_D$ 的归纳假设一起表明,
无论选取 $D$ 的哪个分支 $D'$, 层 $i^*R^ij_*\underline G$ 都是 $T$-整的.
这就证明了引理~5.6.1.

尚需处理曲线. 在引理~5.6.1 的假设下, 设 $x$ 为 $X-U$ 的点,
$\mathfrak G$ 为 $\widehat{\mathcal O}_{x,X}$ 的分式域的 Galois 群,
$\mathfrak J$ 为惯性群. 仍用 $\underline G$ 表示该层定义的 $\mathfrak G$ 的表示,
并设 $\bar x$ 是 $X$ 的一个位于 $x$ 处的几何点. 有 Galois $\mathfrak G/\mathfrak J$ 模同构
\[
(j_*\underline G)_{\bar x}\simeq\underline G^{\mathfrak J},
\]
\[
(R^1j_*\underline G)_{\bar x}\simeq\underline G_{\mathfrak J}(-1),
\]
及
\[
(R^ij_*\underline G)_{\bar x}=0\qquad\text{当 }i\geq2.
\]
断言由推论~5.3(a) 和引理~5.3.1 得出.

% Source: scan idx 407 | printed folio 400.
\section*{参考文献}
\begin{enumerate}[label={\textup{[\arabic*]}},leftmargin=2.4em]
\item S.~S. Abhyankar, \emph{Resolution of Singularities of Embedded Algebraic Surfaces}, Academic Press, New York, 1966.
\item J.-P.~Jouanolou, \emph{Cohomologie $\ell$-adique des schémas}, 将刊于 \emph{Lecture Notes in Mathematics}.
\item S.~Lang 与 A.~Weil, \emph{Number of Points of Varieties in Finite Fields}, \emph{Amer. J. Math.} \textbf{76} (1954), 819--827.
\item A.~Weil, \emph{Sur les courbes algébriques et les variétés qui s'en déduisent}, \emph{Actualités Sci. Indust.} no.~1041, Hermann, Paris, 1948.
\end{enumerate}

\medskip
\noindent\emph{底本注.}
直接对照权威底本, 将分解图中的 $j$ 标签恢复到箭头下侧,
并在曲线情形恢复印本的 Fraktur 字体 $\mathfrak G$.
照录法文层保留了印本的单数 \emph{morphisme}、假设 $(*)$ 前的标点冲突,
以及简写引用 6.2(i). 校订法文层和英文层采用复数 \emph{morphismes}、正常句子标点,
以及准确的引用引理 5.6.2(i).
